\appendix[Physical interpretation of the phase function's Hessian]
\label{app:hessian}
This section briefly summarizes the physical interpretation of the phase function's Hessian and formulates directional second derivatives for \eqref{eq:KA} in the present geometry in order to evaluate it by the SPA.
For further elaboration refer to \cite[App. B]{Firtha2019phd}.

Assume the phase of the arbitrary field given by \eqref{eq:general_field}!
The 3x3 Hessian of the phase function $\mH^P(\vx)$ (containing the second derivatives, w.r.t combinations of $x,y,z$) is given by
\begin{equation}
  \mH^P(\vx) = -
  \begin{bmatrix} 
    \phi^{P''}_{xx}(\vx,\omega) & \phi^{P''}_{xy}(\vx,\omega) & \phi^{P''}_{xz}(\vx,\omega) \\[.7em]
    \phi^{P''}_{xy}(\vx,\omega) & \phi^{P''}_{yy}(\vx,\omega) & \phi^{P''}_{yz}(\vx,\omega) \\[.7em]
    \phi^{P''}_{xz}(\vx,\omega) & \phi^{P''}_{yz}(\vx,\omega) & \phi^{P''}_{zz}(\vx,\omega) \\[0.5em]    \end{bmatrix}.
\end{equation}
The matrix can be written in the spectral form of
\begin{equation}
  \mH^P(\vx) = - \underbrace{\begin{bmatrix} \mathbf{v}^P_1 \, |\, \mathbf{v}^P_2 \\\end{bmatrix}}_{\mathbf{V}^P(\vx)} 
  \begin{bmatrix} \kappa^P_1 & 0 \\[.0em] 0 & \kappa^P_2 \\[0.0em] \end{bmatrix}
  \begin{bmatrix} \mathbf{v}^P_1 \, | \, \mathbf{v}^P_2 \\\end{bmatrix}^{\mathrm{T}},
\end{equation}
with 
\begin{itemize}
  \item $\kappa^P_{1,2} = \frac{1}{\rho^P_{1,2}}$ being the maximal and minimal curvatures of the wavefront, and $\rho^P_{1,2}$ being the corresponding principal radii, being the radii of the maximal and minimal osculating circles, tangential with the wavefront at $\vx$.
  \item $\mathbf{v}^P_{1,2}$ are the corresponding tangential vectors at $\vx$, pointing into the direction of the minimal and maximal curvatures, being both perpendicular to $\vhk^P(\vx)$.
  \item The third eigenvector is given by the local wavenumber vector itself, with the corresponding eigenvalue being 0, explaining the dimensionality reduction of the matrix from 3x3 to 2x2.
\end{itemize}
The geometrical properties are illustrated in Figure \ref{Fig:local_wave_curvature}.

In the present problem we require the phase function's Hessian for the integrand of the Kirchhoff approximation \eqref{eq:KA}, given by
{\fontsize{8pt}{\baselineskip}
\selectfont
\begin{multline}
  \mH^{P+G}(\vxo) = \\
  -\mathbf{V}^P(\vxo)
  \underbrace{\begin{bmatrix} \frac{1}{\rho^G_h(\vx-\vxo)} + \frac{1}{\rho^P_h(\vxo)} & 0 \\[.0em] 0 & \frac{1}{\rho^G_v(\vx-\vxo)} + \frac{1}{\rho^P_v(\vxo)} \\[0.0em] \end{bmatrix}}_{\boldsymbol{\kappa}^{P+G}(\vxo)}
  \mathbf{V}^P(\vxo)^{\mathrm{T}}.
\end{multline}}
In this expression it was exploited that propagating by the Green's function does not change the wavefront's tangential vectors $\mathbf{V}^P$, hence $\mathbf{V}^{P+G} = \mathbf{V}^P$ holds. 
Furthermore, we restricted our investigation to sound fields with $\hat{k}^P_z = 0$ in which case the principal curvatures simplify to horizontal and vertical components, with the corresponding tangential vectors of $\mathbf{V}^P(\vxo) = \begin{bmatrix} \mathbf{t}^P(\vxo) & \hat{\mathbf{z}} \\\end{bmatrix}$.
Here $\mathbf{t}^P(\vxo)$ denotes the tangential vector to the wavefront in the horizontal plane and $\hat{\mathbf{z}}$ is the vertical unit vector.

With the above mathematical description we can simply formulate cases when the second derivative into an arbitrary direction $\mathbf{v}_0(\vxo)$ is required by evaluating
\begin{multline}
  \phi^{P+G''}_{v_0v_0} =
-  \mathbf{v}_0(\vxo)^{\mathrm{T}}\left( \mathbf{V}^P(\vxo) \boldsymbol{\kappa}(\vxo)^{P+G} \mathbf{V}^P(\vxo)^{\mathrm{T}} \right)\mathbf{v}_0(\vxo).
\end{multline}

In practical cases 

\paragraph{Application for the vertical SPA}
For the case of the vertical SPA $\mathbf{v}_0(\vxo) = \hat{\mathbf{z}}$ is given.
Since $\mathbf{v}_0(\vxo)^{\mathrm{T}} \mathbf{V}^P(\vxo) = [0,1]$, therefore, the required derivative is simply given as
\begin{equation}
  \phi^{P+G''}_{zz} = -\left( \frac{1}{\rho^G_v(\vx-\vxo)} + \frac{1}{\rho^P_v(\vxo)}\right).
  \label{eq:curvature_vert}
\end{equation}

\paragraph{Application for the horizontal SPA}
%On the other hand, when the second derivative is required in the horizontal plane, in case of an arbitrary shaped SSD contour the SSD curvature has to be taken into account.
Assume that we have an arc length parametrization of the contour given by $s(\vxo)$, with the tangential vector (first derivative of the contour) given by $\mathbf{s}(\vxo)$ and the curvature function (second derivative) by $\kappa_{\mathrm{SSD}}(\vxo) = 1 / \rho_{\mathrm{SSD}}(\vxo)$.  
In this geometry the second derivative along $s(\vxo)$ can be written by applying the chain rule to $\phi(s(\vxo))$ as
\begin{multline}
  \phi_{ss} = \mathbf{s}(\vxo)^{\mathrm{T}} \, \mH(\vx) \, \mathbf{s}(\vxo) + 2\kappa_{\mathrm{SSD}}(\vxo) \left< \Delta \phi(\vxo) \cdot \mathbf{n}(\vxo) \right>
\end{multline}

For the first term by using the spectral decomposition of the Hessian, $\mathbf{s}(\vxo)^{\mathrm{T}} \mathbf{V}^P(\vxo) = [\left<\mathbf{s}(\vxo) \cdot \mathbf{v}_1^P(\vxo)\right>,0]$ holds, i.e.
\begin{multline}
  \small
  \mathbf{s}(\vxo)^{\mathrm{T}} \, \mH(\vx) \, \mathbf{s}(\vxo)  = \\ =
  -\left<\mathbf{s}(\vxo) \cdot \mathbf{v}_1^P(\vxo)\right>^2 \left(\frac{1}{\rho^G_h(\vx-\vxo)} + \frac{1}{\rho^P_h(\vxo)} \right).
\end{multline}
Finally, it can be realized that the inner product can be written in terms of horizontal, perpendicular components (which for $\mathbf{v}_1^P(\vxo)$ is given by $\vhk^P(\vxo)$) the equation takes the form
\begin{multline}
  \phi^{P+G''}_{ss} = 
   -\hat{k}_{\mathrm{n}}(\vxo)^2 \left(\frac{1}{\rho^G_h(\vx-\vxo)} + \frac{1}{\rho^P_h(\vxo)} \right) - \\
   -2 \left( \hat{k}^P_{\mathrm{n}}(\vxo) - \hat{k}^G_{\mathrm{n}}(\vx-\vxo)\right) \kappa_{\mathrm{SSD}}(\vxo).
   \label{eq:curvature_horiz}
\end{multline}
Note that the above equation is equivalent with the result given by \cite[Eq. 8.4.22.]{Bleistein1984}, where the local coordinate system is chosen to coincide with the principal curvatures of the SSD specializing the virtual field to be a point source, while here the SSD is invariant along the $z$-dimension, allowing to fix the coordinate frame to the virtual source.  

